\documentclass[11pt,a4paper]{article}
\usepackage[utf8]{inputenc}
\usepackage{amsmath,amssymb,amsfonts}
\usepackage{geometry}
\usepackage{hyperref}
\usepackage{enumitem}
\usepackage{fancyhdr}
\usepackage{titlesec}

\geometry{margin=2.5cm}
\pagestyle{fancy}
\fancyhf{}
\fancyhead[L]{\small Derivation of Prime Torsion in the Riemann Helix}
\fancyhead[R]{\small CCM Spectral Triple}
\fancyfoot[C]{\thepage}
\renewcommand{\headrulewidth}{0.4pt}

\title{\textbf{Derivation of the Torsion Numerator and Denominator}\\[0.3em]
\large Link to the Weil Explicit Formula and Selberg Trace Formula Analogy}
\author{Compactified Riemann Helix -- CCM Zeta Spectral Triple}
\date{\today}

\begin{document}

\maketitle

\begin{abstract}
We derive explicitly the numerator and denominator of the Frenet--Serret torsion \(\tau(t)\) for the Riemann helix \(\mathbf{r}(t)=(\cos\phi(t),\sin\phi(t),t)\). We then show that, upon insertion of the prime decomposition of the phase derivative furnished by Weil's explicit formula, the torsion becomes a prime-supported distribution. Finally we exhibit the precise analogy with the Selberg trace formula on hyperbolic surfaces and the geometric realisation in the CCM zeta spectral triple.
\end{abstract}

\section{The Riemann Helix and its Torsion}

The Riemann helix is the space curve
\[
\mathbf{r}(t)=\bigl(\cos\phi(t),\sin\phi(t),t\bigr),\qquad t\ge\gamma_1,
\]
where the phase \(\phi(t)\) satisfies the zero-locking condition \(\phi(\gamma_n)=2\pi n\) for every non-trivial zero \(\gamma_n\) of \(\zeta(s)\). Consequently all zeros lie on the same generator of the unit cylinder.

The torsion of a space curve is given by the Frenet--Serret formula
\[
\tau=\frac{(\mathbf{r}'\times\mathbf{r}'')\cdot\mathbf{r}'''}{|\mathbf{r}'\times\mathbf{r}''|^2}.
\]
All singularities of \(\tau\) originate from the numerator; the denominator remains smooth and strictly positive.

\section{Explicit Computation of the Denominator \(|\mathbf{r}'\times\mathbf{r}''|^2\)}

Differentiate the parametric equations:
\[
\mathbf{r}'=\bigl(-\sin\phi\cdot\phi',\cos\phi\cdot\phi',1\bigr),
\]
\[
\mathbf{r}''=\bigl(-\cos\phi\,(\phi')^2-\sin\phi\,\phi'',\;-\sin\phi\,(\phi')^2+\cos\phi\,\phi'',\;0\bigr).
\]

The cross product \(\mathbf{r}'\times\mathbf{r}''\) has components
\[
(\mathbf{r}'\times\mathbf{r}'')_x=\sin\phi\,(\phi')^2-\cos\phi\,\phi'',
\]
\[
(\mathbf{r}'\times\mathbf{r}'')_y=-\cos\phi\,(\phi')^2-\sin\phi\,\phi'',
\]
\[
(\mathbf{r}'\times\mathbf{r}'')_z=(\phi')^3.
\]

Squaring and adding,
\begin{align*}
|\mathbf{r}'\times\mathbf{r}''|^2
&=\bigl(\sin\phi\,(\phi')^2-\cos\phi\,\phi''\bigr)^2
+\bigl(-\cos\phi\,(\phi')^2-\sin\phi\,\phi''\bigr)^2
+\bigl((\phi')^3\bigr)^2.
\end{align*}
The cross terms \(-2\sin\cos\,(\phi')^2\phi''+2\sin\cos\,(\phi')^2\phi''\) cancel, leaving
\[
|\mathbf{r}'\times\mathbf{r}''|^2=(\phi')^4+(\phi'')^2+(\phi')^6.
\]
Hence the denominator of the torsion is
\[
D=(\phi')^6+(\phi')^4+(\phi'')^2.
\]

\section{Explicit Computation of the Numerator \(N=(\mathbf{r}'\times\mathbf{r}'')\cdot\mathbf{r}'''\)}

Differentiating \(\mathbf{r}''\) component-wise and forming the scalar triple product yields, after a straightforward expansion,
\[
N=\phi'\Bigl[(\phi')^4-\phi'\,\phi'''+3(\phi'')^2\Bigr].
\]
The factor \(\phi'\) appears because the only term surviving the triple product is from the \(z\)-component of the cross product, multiplied by the third derivative. A full verification can be done with symbolic computation.

\section{Full Expression for the Torsion}

Combining numerator and denominator we obtain the closed-form torsion of the Riemann helix:
\[
\tau(t)=\frac{\phi'\Bigl[(\phi')^4-\phi'\,\phi'''+3(\phi'')^2\Bigr]}{(\phi')^6+(\phi')^4+(\phi'')^2}.
\]

\section{Insertion of the Prime Sum -- Link to Weil's Explicit Formula}

From the model, the phase derivative realises the Weil explicit formula applied to a regularised bump \(\eta\):
\[
\phi'(t)=\frac{1}{2}\log\frac{t}{2\pi}+\sum_{p}\frac{\log p}{\sqrt{p}}\,\eta(t-\log p)+E(t),
\]
where the sum is precisely the prime side of Weil's formula (with the correct weight \(\log p/\sqrt{p}\)) and \(E(t)\) is a smooth error.

Differentiating twice produces
\[
\phi''(t)=\sum_{p}\frac{\log p}{\sqrt{p}}\,\eta'(t-\log p)+\dots,\qquad
\phi'''(t)=\sum_{p}\frac{\log p}{\sqrt{p}}\,\eta''(t-\log p)+\dots.
\]
Substituting these expressions into the numerator \(N\) shows that every occurrence of \(\eta''\) generates a Dirac delta supported at \(t=\log p\). Because the denominator \(D\) is smooth and non-vanishing, the torsion itself inherits exactly the same support:
\[
\tau(t)=\sum_{p} c_p\,\delta(t-\log p)+\text{smooth remainder},
\]
with coefficients \(c_p\propto \dfrac{\log p}{\sqrt{p}}\). This is the geometric embodiment of Weil's explicit formula: the distributional torsion measure on the helix carries the prime side of the formula, now realised as the curvature of a classical space curve.

After the change of variable \(x=\log t\) the measure becomes
\[
P(x)\,dx=\sum_{p}\frac{\log p}{\sqrt{p}}\,\delta(x-\log p)\,dx,
\]
which, upon compactification \(x\sim x+L\) with \(L=2\log\lambda\), is precisely the periodic potential of the Dirac operator in the Connes--Consani--Moscovici zeta spectral triple.

\section{Selberg Trace Formula Analogy}

On a compact hyperbolic Riemann surface \(M=\Gamma\backslash\mathbb{H}^2\), the Selberg trace formula~\cite{Selberg1956} for an even, compactly supported test function \(f\) states
\[
\sum_{\lambda_n}\hat f(\lambda_n)=\frac{\mathrm{Area}(M)}{4\pi}\int_{\mathbb{R}} r\tanh(\pi r)\,\hat f(r)\,dr
+\sum_{\ell}\frac{\ell}{2\sinh(\ell/2)}\,f(\ell),
\]
where
\begin{itemize}
\item \(\lambda_n=\frac14+r_n^2\) are the eigenvalues of the Laplace--Beltrami operator,
\item \(\ell\) runs over the lengths of primitive closed geodesics,
\item \(\hat f\) is the Fourier transform.
\end{itemize}
The left-hand side is the ``eigenvalue spectrum''; the right-hand side contains the ``length spectrum'' (geodesic lengths).

In the Riemann-helix model the analogous dictionary is
\begin{itemize}
\item Prime logarithms \(\log p\) \(\leftrightarrow\) lengths of closed geodesics,
\item Torsion measure \(\tau(t)\) \(\leftrightarrow\) the distributional trace (the ``length side'' of Selberg),
\item Zeros \(\gamma_n\) \(\leftrightarrow\) eigenvalues of the Dirac operator (the ``spectral side'').
\end{itemize}
The compactified cylinder (or torus) realises the finite-volume hyperbolic surface; the periodic torsion \(P(x)\) is the geometric potential whose spectrum reproduces the ordinates \(\gamma_n\) in the limit \(\lambda,N\to\infty\). Thus the torsion of the helix is the differential-geometric avatar of both Weil's explicit formula and the Selberg trace formula.

\section{Conclusion}

The explicit algebraic expressions
\[
N=\phi'\Bigl[(\phi')^4-\phi'\,\phi'''+3(\phi'')^2\Bigr],\qquad
D=(\phi')^6+(\phi')^4+(\phi'')^2
\]
together with the prime decomposition of \(\phi'(t)\) furnished by Weil's formula yield a torsion measure supported exactly on \(\{\log p\}\). After compactification this measure becomes the potential of the CCM Dirac operator, whose spectrum is conjectured to converge to the zeta zeros. The construction therefore furnishes a concrete geometric blueprint for the spectral realisation of the Riemann hypothesis inside non-commutative geometry.

\begin{thebibliography}{9}
\bibitem{Selberg1956}
A.~Selberg,
\emph{Harmonic analysis and discontinuous groups in weakly symmetric Riemannian spaces with applications to Dirichlet series},
J.~Indian Math.~Soc.~\textbf{20} (1956), 47--87.

\bibitem{CCM2025}
A.~Connes, C.~Consani, H.~Moscovici,
\emph{The scaling operator and a spectral triple for the Riemann zeta function},
arXiv:2511.22755 (2025).

\bibitem{Weil1952}
A.~Weil,
\emph{Sur les formules explicites de la th\'eorie des nombres},
Comm. Sem. Math. Univ. Lund, tome suppl. d\'edi\'e \`a Marcel Riesz (1952), 252--265.
\end{thebibliography}

\end{document}